\uuid{wRXj}
\titre{Pgcd}
\auteur{bodin ; hanani ; mzari}
\organisation{exo7}
\licence{CC BY-NC-SA 4.0}
\datecreate{2020-01-20}
\niveau{L1}
\module{Algèbre}
\chapitre{Polynôme, fraction rationnelle}
\sousChapitre{Division euclidienne}
\difficulte{2}
\theme{Arithmétique des polynômes, pgcd}
\exo7qcmid{lille-1.polynomes.013}

\contenu{
\begin{question}
\qclassification{105.01,105.02}
\qtags{motcle=pgcd}
Soient $A(X) = X^4 - 2 X^3 - 4 X^2 + 2 X + 3$ et
$B(X) = X^4 - 2 X^3 - 3 X^2$ des polynômes de $\Rr[X]$.
Notons $D$ le pgcd de $A$ et $B$.
Quelles sont les affirmations vraies  ?
\begin{answers}
    \bad{$X-1$ divise $D$.}

    \good{$X+1$ divise $D$.}

    \good{$D(X) = (X-3)(X+1)$.}

    \bad{$D(X) = (X-3)(X+1)^2$.}
\end{answers}
\begin{explanations}
$A(X) = (X-3)(X+1)^2(X-1)$, $B(X) = X^2(X-3)(X+1)$,
le pgcd est $D = (X-3)(X+1)$.
\end{explanations}
\end{question}
}
